% GENERATED FILE. Edit canonical lessons, then run npm run generate:books.
% canonical-source-hash: fnv1a64:52b24900af596b1b
% canonical-lessons: ES-C640-veranear, ES-C640-automovil, ES-C640-circulacion, ES-C640-furgoneta, ES-C640-conductor

\chapter{Cómo nos movemos}
\label{ch:es-como-nos-movemos}
\hlchaptermodality{640}

\begin{chapteropening}
\textbf{By the end of this chapter:} \emph{I can talk about how I get around: summer trips, cars and vans, the traffic and the driver.}

Complete the last lesson of chapter 640: I can talk about how I get around: summer trips, cars and vans, the traffic and the driver.
\end{chapteropening}

\section[veranear]{veranear — to spend the summer (on holiday)}

\label{lesson:ES-C640-veranear}



\begin{quote}
Before the new one: say the Spanish for to seem, then the Spanish for to keep.
\end{quote}

\subsection*{You'll want to know: veranear}
\textbf{veranear} — \textquotedblleft{}to spend the summer (on holiday)\textquotedblright{}.

\begin{quote}
Veraneamos en la costa todos los años.
\end{quote}

\begin{quote}
\emph{We spend every summer on the coast.}
\end{quote}

Built on \emph{el verano}. \emph{El veraneo} is the summer holiday.

\begin{cousinweb}
\emph{veranear} is built on \emph{verano} \textquotedblleft{}summer\textquotedblright{}, which is inherited from Vulgar Latin \emph{vērānum (tempus)}, built on Latin \emph{vēr} \textquotedblleft{}spring\textquotedblright{} (EtymDB).
\end{cousinweb}

\begin{grammarlens}[title={What you've built}]
One more word for telling, explaining and describing at B1.
\end{grammarlens}

\subsection*{Your turn}
\begin{itemize}
\raggedright
  \item \emph{Say it:} \emph{veranear}
  \item \emph{Say it:} \emph{veranear}, once more
  \item \emph{Say it:} a sentence of your own with \emph{veranear}
  \item \emph{Recall:} say the Spanish for to seem, then the Spanish for to keep, then say \emph{veranear} again
\end{itemize}

\begin{tcolorbox}[breakable,colback=teal!4,colframe=teal!35!black,title={Before you move on}]
What does veranear mean? (\textquotedblleft{}To spend the summer (on holiday)\textquotedblright{}.) Say it once more.
\end{tcolorbox}
\section[el automóvil]{el automóvil — the car, the automobile}

\label{lesson:ES-C640-automovil}



\begin{quote}
Before the new one: say the Spanish for to keep, then the Spanish for to spend the summer.
\end{quote}

\subsection*{You'll want to know: el automóvil}
\textbf{el automóvil} — \textquotedblleft{}the car, the automobile\textquotedblright{}.

\begin{quote}
Mi abuelo compró su primer automóvil en 1970.
\end{quote}

\begin{quote}
\emph{My grandfather bought his first car in 1970.}
\end{quote}

A formal or old-fashioned word; in speech, \emph{el coche} in Spain, \emph{el carro} or \emph{el auto} in Latin America.

\begin{cousinweb}
\emph{automóvil} came from French \emph{automobile}, a word built from Greek \emph{autós} \textquotedblleft{}self\textquotedblright{} (EtymDB) and Latin \emph{mōbilis} \textquotedblleft{}movable\textquotedblright{}: a carriage that moves by itself.
\end{cousinweb}

\begin{grammarlens}[title={What you've built}]
One more word for telling, explaining and describing at B1.
\end{grammarlens}

\subsection*{Your turn}
\begin{itemize}
\raggedright
  \item \emph{Say it:} \emph{el automóvil}
  \item \emph{Say it:} \emph{el automóvil}, once more
  \item \emph{Say it:} a sentence of your own with \emph{el automóvil}
  \item \emph{Recall:} say the Spanish for to keep, then the Spanish for to spend the summer, then say \emph{el automóvil} again
\end{itemize}

\begin{tcolorbox}[breakable,colback=teal!4,colframe=teal!35!black,title={Before you move on}]
What does el automóvil mean? (\textquotedblleft{}The car, the automobile\textquotedblright{}.) Say it once more.
\end{tcolorbox}
\section[la circulación]{la circulación — the traffic (flow); circulation}

\label{lesson:ES-C640-circulacion}



\begin{quote}
Before the new one: say the Spanish for to spend the summer, then the Spanish for the car.
\end{quote}

\subsection*{You'll want to know: la circulación}
\textbf{la circulación} — \textquotedblleft{}the traffic (flow); circulation\textquotedblright{}.

\begin{quote}
A esta hora hay mucha circulación en el centro.
\end{quote}

\begin{quote}
\emph{At this time there is a lot of traffic in the centre.}
\end{quote}

On the radio: \emph{la circulación es lenta} — traffic is slow. It is also the circulation of the blood.

\begin{cousinweb}
\emph{circulación} is built on \emph{circular} \textquotedblleft{}to move round\textquotedblright{}, a learned borrowing from Late Latin \emph{circulāre} (EtymDB), from \emph{circulus} \textquotedblleft{}a small ring\textquotedblright{}, a little \emph{circus}.
\end{cousinweb}

\begin{grammarlens}[title={What you've built}]
One more word for telling, explaining and describing at B1.
\end{grammarlens}

\subsection*{Your turn}
\begin{itemize}
\raggedright
  \item \emph{Say it:} \emph{la circulación}
  \item \emph{Say it:} \emph{la circulación}, once more
  \item \emph{Say it:} a sentence of your own with \emph{la circulación}
  \item \emph{Recall:} say the Spanish for to spend the summer, then the Spanish for the car, then say \emph{la circulación} again
\end{itemize}

\begin{tcolorbox}[breakable,colback=teal!4,colframe=teal!35!black,title={Before you move on}]
What does la circulación mean? (\textquotedblleft{}The traffic (flow); circulation\textquotedblright{}.) Say it once more.
\end{tcolorbox}
\section[la furgoneta]{la furgoneta — the van}

\label{lesson:ES-C640-furgoneta}



\begin{quote}
Before the new one: say the Spanish for the car, then the Spanish for the traffic.
\end{quote}

\subsection*{You'll want to know: la furgoneta}
\textbf{la furgoneta} — \textquotedblleft{}the van\textquotedblright{}.

\begin{quote}
Llevamos las bicis en la furgoneta.
\end{quote}

\begin{quote}
\emph{We are taking the bikes in the van.}
\end{quote}

A small van; much of Latin America says \emph{la camioneta}.

\begin{cousinweb}
\emph{furgoneta} is a little \emph{furgón} \textquotedblleft{}goods wagon\textquotedblright{}, which came from French \emph{fourgon} (EtymDB). Where the French word comes from is not settled.
\end{cousinweb}

\begin{grammarlens}[title={What you've built}]
One more word for telling, explaining and describing at B1.
\end{grammarlens}

\subsection*{Your turn}
\begin{itemize}
\raggedright
  \item \emph{Say it:} \emph{la furgoneta}
  \item \emph{Say it:} \emph{la furgoneta}, once more
  \item \emph{Say it:} a sentence of your own with \emph{la furgoneta}
  \item \emph{Recall:} say the Spanish for the car, then the Spanish for the traffic, then say \emph{la furgoneta} again
\end{itemize}

\begin{tcolorbox}[breakable,colback=teal!4,colframe=teal!35!black,title={Before you move on}]
What does la furgoneta mean? (\textquotedblleft{}The van\textquotedblright{}.) Say it once more.
\end{tcolorbox}
\section[el conductor]{el conductor — the driver}

\label{lesson:ES-C640-conductor}



\begin{quote}
Before the new one: say the Spanish for the traffic, then the Spanish for the van.
\end{quote}

\subsection*{You'll want to know: el conductor}
\textbf{el conductor} — \textquotedblleft{}the driver\textquotedblright{}.

\begin{quote}
El conductor del autobús es muy amable.
\end{quote}

\begin{quote}
\emph{The bus driver is very kind.}
\end{quote}

A woman is \emph{la conductora}. In Spain the verb is \emph{conducir}.

\begin{cousinweb}
\emph{conductor} is a learned borrowing from Latin \emph{conductor}, which meant \textquotedblleft{}one who hires, a contractor\textquotedblright{} (Lewis \& Short); it is built on \emph{condūcere} \textquotedblleft{}to bring together\textquotedblright{}, from \emph{dūcere} \textquotedblleft{}to lead\textquotedblright{} (EtymDB).
\end{cousinweb}

\begin{grammarlens}[title={What you've built}]
One more word for telling, explaining and describing at B1.
\end{grammarlens}

\subsection*{Your turn}
\begin{itemize}
\raggedright
  \item \emph{Say it:} \emph{el conductor}
  \item \emph{Say it:} \emph{el conductor}, once more
  \item \emph{Say it:} a sentence of your own with \emph{el conductor}
  \item \emph{Recall:} say the Spanish for the traffic, then the Spanish for the van, then say \emph{el conductor} again
\end{itemize}

\begin{tcolorbox}[breakable,colback=teal!4,colframe=teal!35!black,title={Before you move on}]
What does el conductor mean? (\textquotedblleft{}The driver\textquotedblright{}.) Say it once more.
\end{tcolorbox}
